\documentclass[11pt]{article}
\usepackage[margin=25mm]{geometry}
\usepackage{fontspec}
\setmainfont{Latin Modern Roman}
\usepackage{amsmath,amssymb,mathtools}
\usepackage{unicode-math}
\setmathfont{Latin Modern Math}
\usepackage{xurl,hyperref}
\hypersetup{colorlinks=true,urlcolor=blue}
\setcounter{secnumdepth}{0}
\setlength{\parindent}{0pt}
\setlength{\parskip}{5pt}
\setlength{\emergencystretch}{3em}
\newcommand{\sref}[2]{\hyperlink{source:#1:#2}{[S#2]}}
\newcommand{\cataloguescope}[1]{\paragraph{Recorded target.}#1}
\newcommand{\R}{\mathbb R}
\newcommand{\C}{\mathbb C}
\newcommand{\N}{\mathbb N}
\newcommand{\Z}{\mathbb Z}
\newcommand{\Q}{\mathbb Q}
\newcommand{\F}{\mathbb F}
\newcommand{\eps}{\varepsilon}
\begin{document}
% UnsolvedMath ID 30002840; source catalogue code OWR-13673-010
\setcounter{section}{30002839}
\section{Conformal Field Theory Realization of the Asaeda–Haagerup Subfactor}\label{30002840}
{\small UnsolvedMath ID 30002840 · Mathematical Physics}
\subsection{Definitions and mathematical statement}

\paragraph{Shared notation.}$\mathbb N=\{1,2,\ldots\}$; the standard number fields have their usual meanings. All additional conventions are stated locally.


\paragraph{Definitions and conventions.}
A subfactor is an inclusion $N\subset M$ of von Neumann algebras with trivial centers. Its index and standard invariant encode the size and higher relative-commutant structure of the inclusion. In an operator-algebraic conformal field theory, interval or local observable algebras and sectors can produce subfactor inclusions. To say a subfactor arises in conformal field theory requires identifying the inclusion or invariant that the construction realizes. The Asaeda--Haagerup subfactor is the exceptional finite-index subfactor named in the source; its standard invariant is the specific datum whose conformal realization is sought.

\paragraph{Target retained from the canonical source catalogue.}
\paragraph{Catalogue wording retained for traceability.}
Does the Asaeda–Haagerup subfactor arise from a conformal field theory?

\paragraph{Source scope and qualification.}
This is a named special case of the general realization question, IDs 30002833 and 30004278. An unrelated subfactor with similar index does not realize the target standard invariant.
\subsection{Short English statement}
Can conformal field theory realize the exceptional Asaeda--Haagerup subfactor?
\subsection{Research attempt}
\paragraph{Model, date and investigated scope.}
OpenAI GPT-6 Astra (configured model \texttt{gpt-6-astra}), ultra reasoning; generated 11 October 2026. Checks the source Asaeda–Haagerup dimension relation and rules out a direct pointed-only sector ansatz for its specified algebra-object realization.

\paragraph{Sources and literature check.}
The full supplied definition was reviewed. The abstract and HTML introduction of Grossman–Izumi–Snyder, arXiv:1501.07324 [S3], were actually opened; the index, object dimension and displayed fusion relation used below appear there. No full source proof or fresh CFT-realization status survey was performed.

\paragraph{Context and investigated scope.} The full supplied definition was reviewed. The abstract and HTML introduction of Grossman–Izumi–Snyder, arXiv:1501.07324 [S3], were actually opened; the index, object dimension and displayed fusion relation used below appear there. No full source proof or fresh CFT-realization status survey was performed.

\paragraph{Testing a pointed-only extension ansatz.} The supplied Grossman--Izumi--Snyder source [S3], whose HTML introduction was consulted, describes an algebra object $1\oplus X$ yielding the Asaeda--Haagerup index $(5+\sqrt{17})/2$, with $d(X)=(3+\sqrt{17})/2$. Put $d=(3+\sqrt{17})/2$. Direct multiplication gives $d^2=3d+2$. Its stated fusion relation $X\otimes X\cong1\oplus X\oplus Y$ therefore forces
\[
 d(Y)=d^2-d-1=2d+1=4+\sqrt{17}.
\]
This is a quick numerical consistency check on the noninvertible sector required by that construction.

\paragraph{Why the simplest simple-current route fails.} In a unitary pointed tensor category all simple objects are invertible and have positive dimension one: multiplicativity gives $d(g)d(g^{-1})=1$, and unitarity identifies the two dimensions. Any finite direct sum of such objects has integral dimension. It therefore cannot be tensor-equivalent to a category containing the above $X$ with irrational dimension, nor realize the specified algebra object merely as a sum of invertible sectors. Adding only a finite set of simple currents is an insufficient ansatz for that direct realization.

\paragraph{Remaining gap.} The primary source supplies tensor-category and subfactor constructions, not a conformal net construction used here. A CFT can have noninvertible sectors, and Morita/orbifold constructions can change the presentation, so the pointed obstruction is narrow. The first missing step is a positive-energy local conformal theory with the appropriate noninvertible tensor data and a verified realization of the actual Asaeda--Haagerup standard invariant. Matching the index or an unrelated exceptional subfactor is insufficient.

\paragraph{Verification and adversarial audit.}
Checked d²=3d+2, d(Y), integral dimensions of pointed sums and the limited scope under possible Morita/orbifold changes. Model self-audit only; no independent human proof checking. Prompts A and B were applied to this canonical record and its qualified source scope. Checked d²=3d+2, d(Y), integral dimensions of pointed sums and the limited scope under possible Morita/orbifold changes. Retained unresolved status and zero original-target progress; A local conformal theory with the required noninvertible sectors and matching Asaeda–Haagerup standard invariant is not constructed.

Coordinator checked the irrational dimension relation and the narrow pointed-category obstruction for the displayed algebra object. This additional check was model-assisted, not independent human verification.

\paragraph{Tools, process and contributions.}
Web browsing opened the primary arXiv abstract and HTML introduction; Python 3 prepared files. Algebraic calculations and prompt-B self-audit were performed by the generating model; no theorem prover or independent human check. Contribution prompt A was followed by prompt B. GPT-6 Astra with ultra reasoning drafted the research and performed its stated self-audit; any additional coordinator contribution is identified in the verification and audit. A GPT-6-based Codex coordinator assembled the manuscript and checked repository integration; its exact variant and reasoning configuration were not exposed. No human mathematical contribution, independent human proof check or formal proof certification is supplied. Complete derivations and the consulted source references are included here; private reasoning transcripts are not published.

\paragraph{Final status and first remaining gap.}
\textbf{UNRESOLVED}. A local conformal theory with the required noninvertible sectors and matching Asaeda–Haagerup standard invariant is not constructed.

\section{Completion Estimate}
\noindent\textbf{COMPLETION: 0\%}
The calculation is an auxiliary or restricted attempt and establishes no newly verified portion of the full source target. A local conformal theory with the required noninvertible sectors and matching Asaeda–Haagerup standard invariant is not constructed.

\subsection{Sources}
\begin{itemize}
\item[\hypertarget{source:30002840:1}{S1}] ULAM.AI and the UnsolvedMath Contributors, UnsolvedMath, record 30002840 (OWR-13673-010). Catalogue curation under CC BY 4.0; definitions and English exposition added locally.
\url{https://huggingface.co/datasets/ulamai/UnsolvedMath}.
\item[\hypertarget{source:30002840:2}{S2}] Subfactors and Conformal Field Theory (2015), Oberwolfach Reports, DOI 10.4171/owr/2015/16, pp. 849--926. Primary report accessed 2026-10-11; relevant PDF pages 35, 36.
\url{https://doi.org/10.4171/owr/2015/16}.
\item[\hypertarget{source:30002840:3}{S3}] Grossman--Izumi--Snyder, Asaeda--Haagerup categories (catalogue-listed context; not independently consulted in this preparation).
\url{https://arxiv.org/abs/1501.07324}.
\end{itemize}
\label{end:30002840}
\end{document}
